% Generated draft: requires mathematical and source review.
\documentclass[11pt]{article}
\usepackage[margin=25mm]{geometry}
\usepackage{fontspec}
\setmainfont{Latin Modern Roman}
\usepackage{amsmath,amssymb,mathtools,unicode-math}
\setmathfont{Latin Modern Math}
\usepackage{enumitem,needspace,xurl,hyperref}
\setlength{\parindent}{0pt}
\setlength{\parskip}{5pt}
\setlength{\emergencystretch}{3em}
\newcommand{\N}{\mathbb{N}}
\newcommand{\Z}{\mathbb{Z}}
\newcommand{\Q}{\mathbb{Q}}
\newcommand{\R}{\mathbb{R}}
\newcommand{\C}{\mathbb{C}}
\newcommand{\F}{\mathbb{F}}
\newcommand{\A}{\mathbb{A}}
\newcommand{\PP}{\mathbb{P}}
\newcommand{\E}{\mathbb{E}}
\newcommand{\eps}{\varepsilon}
\newcommand{\dd}{\,\mathrm{d}}
\begin{document}
\section*{UnsolvedMath ID: 20001654}
\textbf{Source problem code:} catalogue entry
\subsection{Definitions and mathematical statement}
The source asks whether
\[
 i(T,\mu_{T'})\,i(T',\mu_T)\ge1\qquad(T,T'\in X_n)
\]
and, failing a global answer, where it holds. The phrase ``as above'' refers to missing context: the record does not define \(X_n\), the current map, or its normalization. The aid's interpretation uses Outer space of free simplicial metric \(F_n\)-trees, with finite marked core-graph quotients and \(n\ge2\), and Patterson--Sullivan currents normalized by \(i(T,\mu_T)=1\). A geodesic current is a flip-invariant \(F_n\)-invariant Radon measure on pairs of distinct boundary points. The tree-current pairing extends translation length, \(i(T,\eta_w)=\|w\|_T\) for counting currents. Normalization is essential to the numerical constant one; projectively rescaling a tree with its self-normalized current leaves the cross-product unchanged.

The aid reports a simplexwise result. Fix one finite marked core graph \(\mathcal G\), and consider two strictly positive edge-length vectors on it, with trees in the same open Outer-space simplex. Let \(h\) denote volume entropy, the exponential growth rate for distances in the universal-cover metric. It has scaling \(h(c x)=h(x)/c\). Choose entropy-one representatives \(x,y\), so \(h(x)=h(y)=1\), and define directed Bregman divergence on the positive length cone by
\[
 D_h(y\Vert x)=h(y)-h(x)-dh_x(y-x).
\]
The aid reports the exact factorization for the normalized cross-product
\[
 I(x,y):=i(T_x,\mu_{T_y})i(T_y,\mu_{T_x})
  =(1+D_h(y\Vert x))(1+D_h(x\Vert y)).
\]
It reports \(I(x,y)\ge1\), with equality exactly for the same projective metric on this fixed marked graph, equivalently \(x=y\) among entropy-one representatives. No claim is made for zero edge lengths or different marked-graph simplices.

It further identifies one half of the logarithmic diagonal Hessian with the established entropy metric on the simplex. To state the normalization explicitly, on the entropy-one hypersurface put \(g_{h,x}(v,v)=D^2h_x(v,v)\) for tangent \(v\) with \(dh_x(v)=0\). The reported identity is
\[
 \frac12\left.\frac{d^2}{dt^2}\right|_{t=0}
  \log I(x,y(t))=g_{h,x}(v,v),
 \qquad y(0)=x,\ y'(0)=v,\ h(y(t))=1.
\]
Equivalently the second-variable diagonal Hessian is twice this entropy metric. The exact conventions of the cited established metric and Patterson--Sullivan assignment must be checked, since their defining references are not supplied. A local Hessian identity does not itself construct a global Weil--Petersson metric or settle positivity across simplex changes.

The source motivation is a WP-type metric on \(X_n\), perhaps on an \(\operatorname{Out}(F_n)\)-invariant subspace or manifold. Kasra Rafi discusses Teichm\"uller-space currents and Bonahon's metric, and warns that an analogous quadratic-differential/Lipschitz construction fails positivity. That warning concerns the comparison construction, not a supplied counterexample to the displayed Outer-space inequality. The text then introduces a following question from Martin Lustig with a second part added by ``Arnaud,'' mentioning an uncertain ``Huber[?]'' and Besson attribution for decompositions of \(\R\)-trees and dual laminations. That adjacent context lacks the actual next question and is not folded into the present inequality. The aid reports the global cross-simplex question unresolved; all partial formulas are attributed unreviewed aid claims.
\subsection{Short English statement}
With explicit self-normalization, the aid reports the Patterson--Sullivan product inequality on each open marked-graph simplex, an exact Bregman factorization, and a Hessian equal to twice the entropy metric. The global inequality and a global WP-type construction remain unresolved in the supplied account.
\subsection{Sources}
\begin{itemize}
\item Local UnsolvedMath registry, permanent ID \(20001654\), source code \texttt{catalogue entry}; \texttt{aim-geometric-group-theory-notes.json}, \texttt{notes[162]}, tag \texttt{problem}.
\item AIM workshop ``The geometry of the outer automorphism group of a free group,'' 2010, Problem 5.1: \url{https://aimath.org/WWN/outerauto/outerauto.pdf}. Source type PDF; statement status \texttt{reconstructed\_unverified}. The supplied extracted text was used; no PDF was accessed or compiled.
\item The aid names Hilion's symmetric Patterson--Sullivan product and the established entropy metric without authoritative full citations. Original discussion names Kasra Rafi, Bonahon, Martin Lustig, ``Arnaud,'' uncertain ``Huber[?],'' and Besson. The absent preceding definitions and following question are not reconstructed. Original literature field is empty.
\item Research attempt 1: \texttt{partial\_result}, status at run \texttt{open}, rigor label \texttt{proved}, novelty confidence low; registry status \texttt{partially\_solved}, classification \texttt{PARTIAL-PROGRESS}. Created and updated 14 August 2026. Machine-generated aid, not peer reviewed; independently verify before citation.
\item Ulam AI / UnsolvedMath Contributors; CC BY 4.0.
\end{itemize}
\end{document}
